\section{约束运行时以换取 AI 自主性}

\subsection{``生成任意代码'' 的幻觉}

AI 编程领域的早期炒作（以及当前的部分宣传）假设 LLM 将赋予我们\textbf{无限的运行时灵活性}------用任何语言、任何模式、不受约束地生成代码，AI 自己会搞定一切。

然而，OpenAI 团队的 harness 实践揭示了一个相反的事实：要实现可信赖的、大规模 AI 生成代码的可维护性，\textbf{必须有所取舍}。具体来说，提升信任度和可靠性的方式是\textbf{约束解决方案空间}------通过特定的架构模式、强制的边界、标准化的结构来限制 AI 的自由度。

\begin{importantbox}{灵活性与可维护性的权衡}
核心洞察：为了获得可信赖的 AI 生成代码，我们需要用 prompt、规则和充满技术细节的 harness 来\textbf{放弃}``生成任意代码''的灵活性。这意味着：
\begin{itemize}[leftmargin=1.5em]
    \item 更少的技术选择自由
    \item 更严格的架构约定
    \item 更详细的规范文档
    \item 但换来更高的代码质量和可维护性
\end{itemize}
这是一个经典的工程权衡：约束产生秩序，秩序产生信任。
\end{importantbox}

\subsection{约束作为赋能}

这种``约束带来自主性''的思路在软件工程中并不新鲜。容器化（Docker/Kubernetes）通过约束运行环境来实现部署自由；微服务通过约束模块边界来实现独立开发和部署；类型系统通过约束数据表达来减少运行时错误。

在 AI 编程语境中，harness 的约束同样是一种\textbf{赋能}------通过限制 AI 的选择空间，反而让它在受限范围内表现得更好、更可靠。

\subsection{本章小结}

``生成任意代码''的无限灵活性是一个幻觉。实践证明，约束解决方案空间是获得可信赖 AI 生成代码的必要条件。这并非 AI 的局限，而是工程学的普遍规律------有意义的约束是高质量输出的前提。

\newpage

